\section*{Hoofdstuk \ref{ch:b2:current-potential-curves} --- Stroom-potentiaalkrommen}

\begin{solution}{exo:b2:current-potential-curves:1}
Zink wordt geoxideerd: \omterm{def:b2:current-potential-curves:ie-curve}{anodische stroom}, positief. Koperionen worden gereduceerd aan de
koperelektrode: \omterm{def:b2:current-potential-curves:ie-curve}{kathodische stroom}, negatief. De twee stromen zijn even groot in
absolute waarde, omdat dezelfde lading door de keten stroomt.
\end{solution}

\begin{solution}{exo:b2:current-potential-curves:2}
$|i_{\lim}| = 96\,485 \times 0.50 \times 10^{-4} \times 1.6 \times 10^{-9} \times
2.0/(25 \times 10^{-6}) = \qty{6.2e-4}{A}$, dus \qty{0.62}{mA}.
\end{solution}

\begin{solution}{exo:b2:current-potential-curves:3}
Een \omterm{def:b2:current-potential-curves:fast-slow}{snel systeem} snijdt de potentiaalas steil bij zijn nernstpotentiaal; een \omterm{def:b2:current-potential-curves:fast-slow}{traag
systeem} vertoont eromheen een vlak stuk zonder stroom, en zijn takken beginnen
pas na een \omterm{def:b2:current-potential-curves:overpotential}{overspanning}. Op platina: \ce{Fe^3+}/\ce{Fe^2+} (snel);
\ce{O2}/\ce{H2O} (traag).
\end{solution}

\begin{solution}{exo:b2:current-potential-curves:4}
Bij de halvegolfpotentiaal, gelijk aan $E^\circ = \qty{0.77}{V}$ wanneer de twee vormen
even snel diffunderen.
\end{solution}

\begin{solution}{exo:b2:current-potential-curves:5}
Rond \qty{0.77}{V} een kathodische golf van \ce{Fe^3+} (snel), die een plateau bereikt
evenredig met zijn concentratie. Rond \qty{0.34}{V} een tweede golf, waarbij koper
wordt afgezet; zijn plateau komt bij het eerste, ongeveer tweemaal zo hoog bij een gelijke
concentratie en vergelijkbare diffusiecoëfficiënten, omdat $n = 2$. Dan, rond
\qty{0}{V} bij pH 0 op platina (lager zodra de elektrode met koper bedekt is),
de steile muur van de reductie van \ce{H+}.
\end{solution}

\begin{solution}{exo:b2:current-potential-curves:6}
Bij \qty{0.50}{V}: alleen \ce{Fe^3+} wordt gereduceerd tot \ce{Fe^2+}, bij zijn \omterm{def:b2:current-potential-curves:limiting-current}{grensstroom}.
Bij \qty{0.10}{V}: \ce{Fe^3+} wordt gereduceerd en tegelijk wordt koper afgezet,
elk bij zijn \omterm{def:b2:current-potential-curves:limiting-current}{grensstroom}; het totaal is de som van de twee plateaus.
\end{solution}

\begin{solution}{exo:b2:current-potential-curves:7}
pH 0: \qty{0}{V} en \qty{1.23}{V}; pH 14: \qty{-0.83}{V} en \qty{0.40}{V}. Op
platina behoudt het venster zijn breedte en schuift het \qty{59}{mV} per pH-eenheid omlaag,
waarbij elke muur zijn koppel volgt, de anodische nog steeds omhooggeschoven door de
zuurstofoverspanning.
\end{solution}

\begin{solution}{exo:b2:current-potential-curves:8}
Bij \qty{1.0}{V} wordt \ce{Fe^2+} geoxideerd op zijn plateau (\omterm{def:b2:current-potential-curves:ie-curve}{anodische stroom}
evenredig met $[\ce{Fe^2+}]$) en wordt eventueel \ce{Ce^4+} gereduceerd op zijn plateau
(\omterm{def:b2:current-potential-curves:ie-curve}{kathodische stroom} evenredig met $[\ce{Ce^4+}]$). Vóór de equivalentie daalt de
\omterm{def:b2:current-potential-curves:ie-curve}{anodische stroom} lineair met het toegevoegde volume; daarna groeit de \omterm{def:b2:current-potential-curves:ie-curve}{kathodische
stroom} lineair. De twee rechten snijden de nul bij de equivalentie.
\end{solution}

\begin{solution}{exo:b2:current-potential-curves:9}
De plateaus verdubbelen ($|i_{\lim}| \propto 1/\delta$). De halvegolfpotentiaal
verandert niet, omdat de laag voor beide vormen dezelfde is; evenmin de
nulstroompotentiaal van een \omterm{def:b2:current-potential-curves:fast-slow}{snel systeem}, die de nernstpotentiaal is.
\end{solution}

\begin{solution}{exo:b2:current-potential-curves:10}
Zoals in het hoofdstuk met $i_{\lim,c} = 0$: $E = E_{1/2} + (RT/nF)\ln[i/(i_{\lim} -
i)]$. Met briggse logaritmen is $E = E_{1/2} + (RT\ln 10/nF)\log[i/(i_{\lim} - i)]$,
een rechte met helling $0.0592/n$ volt bij \qty{298}{K}, die de
$E$-as snijdt bij $E_{1/2}$.
\end{solution}

\begin{solution}{exo:b2:current-potential-curves:11}
Op zichzelf neemt het zink de potentiaal aan waarbij zijn \omterm{def:b2:current-potential-curves:ie-curve}{anodische stroom} (snelle
oxidatie van zink) gelijk is aan de \omterm{def:b2:current-potential-curves:ie-curve}{kathodische stroom} van de reductie van \ce{H+} op zink,
die klein is omdat die reductie op zink traag is: het oplossen verloopt
traag. In contact met platina kunnen de elektronen die het zink afgeeft \ce{H+}
reduceren op het platina, waar de reductie snel is: de \omterm{def:b2:current-potential-curves:ie-curve}{kathodische stroom} bij een gegeven
potentiaal is veel groter, het evenwicht wordt bij een grotere stroom bereikt, en het
zink lost sneller op terwijl er waterstof ontstaat op het platina.
\end{solution}

\begin{solution}{exo:b2:current-potential-curves:12}
Op platina begint de oxidatie van water bij pH 0 rond \qty{1.6}{V}: vóór
\qty{2.0}{V} zou de stroom opgaan aan het maken van dizuurstof. Een elektrode waarop
de oxidatie van water een grotere \omterm{def:b2:current-potential-curves:overpotential}{overspanning} vereist, duwt de muur voorbij
\qty{2.0}{V}, zodat het deeltje binnen het venster geoxideerd kan worden.
\end{solution}

\begin{solution}{pb:b2:current-potential-curves:1}
\textbf{1.} \ce{O2 + 2H2O + 4e- -> 4OH-}.
\textbf{2.} \ce{Ag + Cl- -> AgCl + e-} (vier keer); totaal \ce{O2 + 2H2O + 4Ag +
4Cl- -> 4AgCl + 4OH-}.
\textbf{3.} $E = 1.229 - 0.0592 \times 7 + (0.0592/4)\log 0.21 = \qty{0.80}{V}$.
\textbf{4.} $E^\circ(\ce{AgCl}/\ce{Ag}) = \qty{0.22}{V}$ met chloride van activiteit 1:
$-0.60 + 0.22 = \qty{-0.38}{V}$ op de waterstofschaal.
\textbf{5.} De reductie van \ce{O2} is een \omterm{def:b2:current-potential-curves:fast-slow}{traag systeem}: onder haar nernstpotentiaal
is een \omterm{def:b2:current-potential-curves:overpotential}{overspanning} van enkele tienden volt nodig voordat de stroom
groeit, en meer om het plateau te bereiken.
\textbf{6.} Op het plateau hangt de stroom alleen af van de aanvoer van \ce{O2},
niet van kleine verschuivingen van de potentiaal: hij meet de concentratie.
\textbf{7.} De reductie van water tot diwaterstof (de kathodische muur), die
een stroom zou toevoegen die niets met zuurstof te maken heeft.
\textbf{8.} Lineair over het membraan, van $c$ aan de kant van het monster tot nul aan
de kathode.
\textbf{9.} $j = D_m c/L$ per eenheid oppervlakte.
\textbf{10.} $i = -4FAD_m c/L$ (kathodisch), $|i| = 4FAD_mc/L$.
\textbf{11.} Het membraan is de traagste stap: diffusie erdoorheen is veel
trager dan in het water erbuiten, dus het membraan speelt de rol van de
\omterm{def:b2:current-potential-curves:limiting-current}{diffusielaag}.
\textbf{12.} De concentratiegradiënt ligt in het membraan, waarvan de dikte
niet van de roering afhangt (zolang het water erbuiten ververst wordt);
een afzetting op het membraan voegt dikte toe en verlaagt de stroom.
\textbf{13.} $0.2095 \times (101.325 - 3.17) = \qty{20.56}{kPa}$.
\textbf{14.} $c = 1.3 \times 10^{-5} \times 20\,560 = \qty{0.27}{mol/m^3}$ ($0.267$).
\textbf{15.} $0.267 \times 32.00 = \qty{8.6}{mg/L}$ (\unit{g/m^3}).
\textbf{16.} $|i| = 4 \times 96\,485 \times 1.0 \times 10^{-5} \times 1.0 \times 10^{-10}
\times 0.267/(25 \times 10^{-6}) = \qty{4.1}{\micro A}$.
\textbf{17.} $4.05/8.55 = \qty{0.474}{\micro A}$ per \unit{mg/L}.
\textbf{18.} De membraandikte en haar diffusiecoëfficiënt zijn niet nauwkeurig
bekend en veranderen met de ouderdom en de temperatuur; een meting in een bekende
oplossing omvat ze allemaal.
\textbf{19.} Nul (in de praktijk een kleine reststroom, die wordt afgetrokken).
\textbf{20.} $2.80/0.474 = \qty{5.9}{mg/L}$.
\textbf{21.} $2.80/4.05 = 69~\%$ van de verzadiging.
\textbf{22.} Zuurstof die verbruikt wordt door de afbraak van organisch materiaal (rioolwater, dode planten),
sneller dan hij wordt aangevoerd door de lucht en door fotosynthese.
\textbf{23.} $2.80 \times 10^{-6}/(4 \times 96\,485) \times 3600 = \qty{2.6e-8}{mol}$
per uur: verwaarloosbaar naast de zuurstof van zelfs één liter water
(\qty{1.8e-4}{mol}).
\textbf{24.} Zowel de oplosbaarheid van \ce{O2} (de \omterm{prop:b2:chemical-potential:henry}{henryconstante}) als de diffusie
door het membraan verandert met de temperatuur: de gevoeligheid die bij
\qty{25}{\degreeCelsius} gevonden werd, geldt niet meer.
\textbf{25.} Het monster bevat $\boldsymbol{\approx \qty{5.9}{mg/L}}$ opgeloste
zuurstof.
\end{solution}
